% Figure 0.1 --- the one reader's map the Book keeps: five archetypes, drawn as
% flowing paths of varying width through the Book's four parts, replacing the
% six per-chapter Reader's Map sections and abstracts this edition used to
% carry (and promise standalone editions would keep). Full-width: this is the
% one figure important enough to claim the margin column too.
%
% Drawn at its native size (about 7in wide) and left otherwise unwrapped: the
% preamble's own overflow safety net (coordination-papers-mega-volume-preamble.tex,
% "Overflow safety net", \pgfsys@typesetpicturebox) already intercepts every
% tikzpicture wider than \linewidth, sets it across the column and the margin
% at \pdfullwidth (6in), and rescales it with \resizebox itself if it is wider
% still -- that is the ONE way a full-width figure claims the margin column in
% this book, not a hand-rolled adjustwidth+resizebox. A hand-rolled wrapper
% was tried first here and produced a figure that ran off the right edge of
% the page: the safety net's own \resizebox, downstream of a \makebox[\textwidth]
% whose declared width is only \textwidth regardless of the margin content
% already sitting past it, re-scaled that already-correct result back up.
% Removing the manual wrapper and trusting the existing hook fixed it.
%
% Encoding, deliberately: position (which lane) and width/line-style (fat solid
% = read closely; medium solid = read; thin dashed = skim) carry the archetype
% distinction, and every path is labelled directly on the page -- no legend.
% Colour is used only for the four PART headers (an axis label, exactly the
% \pd@chaptercolor already used for the margin glyph elsewhere in this book),
% never to distinguish one archetype's path from another's: the Book's seven
% story colours have measured pairs below the legibility floor for full-colour
% readers (a separate branch, claude/figure-palette-separation, is re-stepping
% them), so a figure that leaned on hue to tell two paths apart would have to
% be redrawn once that lands. This one does not depend on it either way.
\begin{figure}[H]
\centering
\begin{tikzpicture}[pd figure]
  % ---- column geometry: four parts, touching, no gaps -----------------------
  \def\bA{0}\def\bB{3.65}\def\bC{7.30}\def\bD{10.95}\def\bE{14.60}
  \def\cA{1.825}\def\cB{5.475}\def\cC{9.125}\def\cD{12.775}
  \def\yTop{10.55}\def\yHead{10.30}\def\ySub{9.75}\def\yBase{0.55}

  % ---- part dividers and header band -----------------------------------
  \foreach \x in {\bB,\bC,\bD} \draw[pd hairline] (\x,\yBase) -- (\x,\yTop);
  \draw[pd hairline] (\bA,\yTop) -- (\bE,\yTop);
  \node[pd panel title,text=pdcobalt,align=center] at (\cA,\yHead) {Part I\\[1pt]\footnotesize\normalfont Ground Truth};
  \node[pd panel title,text=pdteal,align=center] at (\cB,\yHead) {Part II\\[1pt]\footnotesize\normalfont The Cost of Seeing};
  \node[pd panel title,text=pdviolet,align=center] at (\cC,\yHead) {Part III\\[1pt]\footnotesize\normalfont What Survives the Restart};
  \node[pd panel title,text=pdgold,align=center] at (\cD,\yHead) {Part IV\\[1pt]\footnotesize\normalfont Trade Between Strangers};
  \node[pd note] at (\cA,\ySub) {1 \textperiodcentered\ 2 \textperiodcentered\ 3};
  \node[pd note] at (\cB,\ySub) {4};
  \node[pd note] at (\cC,\ySub) {5};
  \node[pd note] at (\cD,\ySub) {6 \textperiodcentered\ 7 \textperiodcentered\ 8};

  % ---- lane geometry: five archetypes, top to bottom -------------------
  \def\yPrac{8.55}\def\yEng{6.90}\def\ySec{5.25}\def\yEcon{3.60}\def\yPhil{1.95}
  \def\wFat{9.0}\def\wMed{5.0}\def\wThin{2.0}

  % A river segment: #1 x start, #2 x end, #3 y, #4 line width, #5 extra style
  \newcommand{\rmseg}[5]{\draw[pd rule,line cap=round,line width=#4pt,#5] (#1,#3) -- (#2,#3);}
  % A direct label naming the path, set once, at the far left of Part I.
  \newcommand{\rmlabel}[2]{\node[pd direct label,anchor=east,align=right,text width=3.05cm,font=\bfseries\footnotesize] at (-0.30,#1) {#2};}
  % A short note pinned just above a fat segment, naming what it is fat for.
  \newcommand{\rmnote}[3]{\node[pd note,anchor=south,font=\scriptsize\itshape] at (#1,#2+0.28) {#3};}

  % 1. The Practitioner (buyer, solo developer, founder, newcomer) ---------
  \rmlabel{\yPrac}{THE PRACTITIONER}
  \rmseg{\bA}{\bB}{\yPrac}{\wThin}{dashed}
  \rmseg{\bB}{\bC}{\yPrac}{\wFat}{}
  \rmseg{\bC}{\bD}{\yPrac}{\wFat}{}
  \rmseg{\bD}{\bE}{\yPrac}{\wMed}{-{Stealth[length=5mm,width=4mm]}}
  \rmnote{\cB}{\yPrac}{digest-with-zoom, consent}
  \rmnote{\cC}{\yPrac}{role vs.\ person}

  % 2. The Systems Engineer (builds and verifies the mechanism) -----------
  \rmlabel{\yEng}{THE SYSTEMS ENGINEER}
  \rmseg{\bA}{\bB}{\yEng}{\wFat}{}
  \rmseg{\bB}{\bC}{\yEng}{\wFat}{}
  \rmseg{\bC}{\bD}{\yEng}{\wMed}{}
  \rmseg{\bD}{\bE}{\yEng}{\wFat}{-{Stealth[length=5mm,width=4mm]}}
  \rmnote{\cA}{\yEng}{single-writer proof, attenuation}
  \rmnote{\cD}{\yEng}{transfer ceremony, settlement}

  % 3. The Security / Formal-Methods Reviewer ------------------------------
  \rmlabel{\ySec}{THE SECURITY REVIEWER}
  \rmseg{\bA}{\bB}{\ySec}{\wFat}{}
  \rmseg{\bB}{\bC}{\ySec}{\wThin}{dashed}
  \rmseg{\bC}{\bD}{\ySec}{\wThin}{dashed}
  \rmseg{\bD}{\bE}{\ySec}{\wFat}{-{Stealth[length=5mm,width=4mm]}}
  \rmnote{\cA}{\ySec}{threat model, deontic split}
  \rmnote{\cD}{\ySec}{ProVerif / TLA+ status}

  % 4. The Mechanism Designer / Economist ----------------------------------
  \rmlabel{\yEcon}{THE MECHANISM DESIGNER}
  \rmseg{\bA}{\bB}{\yEcon}{\wThin}{dashed}
  \rmseg{\bB}{\bC}{\yEcon}{\wMed}{}
  \rmseg{\bC}{\bD}{\yEcon}{\wMed}{}
  \rmseg{\bD}{\bE}{\yEcon}{\wFat}{-{Stealth[length=5mm,width=4mm]}}
  \rmnote{\cD}{\yEcon}{three-sided market, conservation}

  % 5. The Institutional Theorist (legitimacy, personhood, morality) ------
  \rmlabel{\yPhil}{THE INSTITUTIONAL THEORIST}
  \rmseg{\bA}{\bB}{\yPhil}{\wThin}{dashed}
  \rmseg{\bB}{\bC}{\yPhil}{\wMed}{}
  \rmseg{\bC}{\bD}{\yPhil}{\wFat}{}
  \rmseg{\bD}{\bE}{\yPhil}{\wMed}{-{Stealth[length=5mm,width=4mm]}}
  \rmnote{\cC}{\yPhil}{Locke, Parfit, personhood}
\end{tikzpicture}
\caption{\textbf{One reader's map for the whole book.} Five archetypes, each a
single flowing path left to right through Parts I--IV; line width is the
encoding that matters --- fat where that reader should read closely, thin and
dashed where the part is safe to skim without losing the argument --- and
every path is named on the page rather than in a legend. Position and width
carry the distinction; colour marks only the four parts themselves (the same
per-part hue used throughout this edition), never one archetype against
another, so the figure does not depend on the categorical palette currently
being re-stepped for full-colour legibility (\texttt{claude/figure-palette-separation}).
A sixth reader, the instructor checking exercises, is deliberately not drawn
as a sixth lane: every chapter's own closing \textsc{Exercises} section is
that reader's route, in every part at once, and a sixth flat line spanning
the full width would say nothing a reader could not already infer. Five
paths is close to the legibility ceiling for one figure; a design that
needed a sixth or seventh would be better split into two figures than drawn
here.}
\label{fig:book-reader-map}
\end{figure}
